\section{Method}
\label{sec:method}

\begin{figure*}[tb]
    \centering
    \includegraphics[width=1.0\linewidth]{fig/method_phase5_v1.png}
    \caption{\textbf{Overview of VISTA-Skill.} A frozen VLM executes the active procedural Skill while \red{a sparse, provenance-aware predicate ledger} summarizes episode-local evidence. Before each primitive action, the current compiled Skill produces an expected transition; after execution, an information-isolated branch extracts an evidence-supported transition \red{from public feedback and, only when needed, pre/post RGB observations}. Their typed mismatch is routed to \red{belief refresh, Skill update, or abstention}. Only identifiable, recurrent Skill evidence can produce a one-field candidate, which is promoted after transition checks and \red{paired affected/protected validation}.}
    \label{fig:method_overview}
\end{figure*}

\subsection{Problem Setup and Interface-Only Initialization}

We model embodied execution as a partially observable decision process $\mathcal{M}=(\mathcal{S},\mathcal{A},\mathcal{T},\mathcal{R},\mathcal{X},\mathcal{O})$. At step $t$, the environment occupies an unobserved state $s_t\in\mathcal{S}$. Given an instruction $I$, a first-person observation $x_t\in\mathcal{X}$, a \red{compact episode ledger $B_t$}, and the active Skill $\omega^{(v)}$, a frozen VLM executor produces a plan that is executed as primitive actions,
\begin{equation}
    a_t\sim\pi_{\theta}
    \left(a\mid I,x_t,B_t,\operatorname{render}(\omega^{(v)}),\red{Q_t}\right),
\end{equation}
\red{Here, $Q_t$ contains transient execution emphasis and active recovery obligations.} The executor parameters $\theta$ remain fixed during both Skill acquisition and evaluation.

{\color{red}The environment contract $\mathcal{A}_0$ specifies legal action names, argument formats, and observations, but the primary setting does not initialize the Skill with task strategies, primitive effects, constraints, or recovery rules. Instead, $\omega^{(0)}$ is an \emph{interface-only} Skill: it has five typed but empty fields and no active prediction or termination rule. Consequently, persistent content can be traced to the agent's own interaction evidence rather than to a hand-written seed. All compared evolution methods receive the same initial artifact and update budget; richer hand-written initialization is used only as an initialization-sensitivity control.}

For every primitive action, VISTA-Skill records the pre-action belief and RGB frame, the executed action, the post-action frame, public environment feedback $f_t$, and the active Skill version. \red{The expected transition is materialized before the environment step, whereas the evidence transition is constructed only after $x_{t+1}$ and $f_t$ exist.} The resulting action-level event is
\begin{equation}
    e_t=\left(I,B_t,\omega^{(v)},x_t,a_t,x_{t+1},f_t,
    \widehat{\Delta B}_t,\Delta B_t^{\mathrm{ev}},E_t\right).
\end{equation}
\red{This ordering is part of the information-isolation contract: post-action evidence cannot alter what the Skill predicted before execution.}

\subsection{Structured Skill and Sparse Evidence State}

A versioned procedural Skill contains five semantic fields \red{and synchronized executable views},
\begin{equation}
    \omega^{(v)}=\left\langle
    A^{(v)},P^{(v)},E^{(v)},T^{(v)},C^{(v)};
    \red{R_{\mathrm{pred}}^{(v)},R_{\mathrm{temp}}^{(v)},q_{\mathrm{term}}^{(v)}}
    \right\rangle .
\end{equation}
$A$, $P$, $E$, $T$, and $C$ are the activation, procedure, effect, termination, and constraint fields. Their compact natural-language rendering conditions the executor. \red{The compiled prediction rules $R_{\mathrm{pred}}$ bind action arguments to typed predicates, $R_{\mathrm{temp}}$ represents bounded recovery obligations, and $q_{\mathrm{term}}$ is a typed termination policy. A revision updates the textual field and its compiled representation in the same Skill version, preventing execution and attribution from silently using different rules. Activation is evaluated at trajectory level; procedure, effect, and constraint rules can additionally be compiled against individual actions.} Termination depends only on grounded goal evidence.

{\color{red}Rather than maintaining a dense scene graph, VISTA-Skill stores only predicates relevant to the instruction, active Skill, and recent actions. Each ledger entry has the form
\begin{equation}
    r=\left(\rho,y,c,q,\mathcal{I},\tau,\nu,\kappa,\eta\right),
\end{equation}
where $\rho$ is an instance-grounded predicate, $c$ is confidence, $q$ is the source type, $\mathcal{I}$ contains evidence identifiers, $\tau$ and $\nu$ denote time and view, and $\kappa$ and $\eta$ encode coverage and task relevance. Provenance-aware merging and expiry keep $B_t$ sparse, preserve instance digits, and prevent stale evidence from overwriting newer evidence without provenance.} Predicates retain the original three-valued semantics $y\in\{\mathrm{true},\mathrm{false},\mathrm{unknown}\}$: a temporarily invisible object remains unknown rather than becoming false. Only accepted predicate evidence may mutate $B_t$; expected changes and teacher decisions cannot do so.

{\color{red}Compiled Skill rules also affect execution. Evidence-backed precondition checks establish whether the executor followed an applicable rule, while an action-admission guard enforces active compiled recovery obligations before an environment action is executed. For example, a recurrent failed pick with $\mathrm{near}(o)=\mathrm{false}$ can yield the reusable obligation
\begin{equation}
    \mathbf{G}\!\left(
    \operatorname{fail\mbox{-}pick}(o)\rightarrow
    \left(\neg\operatorname{pick}(o)\ \mathbf{U}\
    \operatorname{recover}(o)\right)
    \right),
\end{equation}
which blocks another pick of the same object from an unchanged state until successful navigation or contrary predicate evidence releases the obligation. Active obligations are also rendered to the executor; guard blocks are logged separately from environment-invalid actions.}

\subsection{Evidence-Decoupled Transition Modeling}

The expected branch is a deterministic compiler operating on the pre-action snapshot,
\begin{equation}
    \widehat{\Delta B}_t=
    P_{\mathrm{exp}}\left(B_t,a_t,\omega^{(v)},\mathcal{A}_0\right).
\end{equation}
\red{In the interface-only discovery setting, $\mathcal{A}_0$ supplies action typing and argument binding but no prior primitive-effect rule; thus the branch predicts only effects, constraints, and termination semantics already accumulated in $\omega^{(v)}$.} Each expected predicate change cites its Skill version, compiled rule, and source field. \red{The deterministic compiler adds no per-action model call.}

The evidence branch receives a deliberately smaller request,
\begin{equation}
    \left(\Delta B_t^{\mathrm{ev}},E_t\right)=
    U_{\mathrm{ev}}\left(I,B_t,a_t,x_t,x_{t+1},f_t\right).
\end{equation}
Its request schema contains neither the active Skill nor $\widehat{\Delta B}_t$, a mismatch, an attribution result, a patch, or a validation outcome. \red{It first parses action-local public feedback and adjacent belief state at no model cost. A visual pre/post provider is triggered only when those sources do not resolve a registered, goal-relevant transition. An authority-aware reliability guard then enforces confidence and coverage thresholds, restricts visual-only claims to predicates local to the executed action, retains explicit action feedback when it conflicts with vision, and downgrades unresolved conflicts to unknown. Goal completion is derived only from independently grounded goal predicates.} The accepted packet is merged into the episode belief,
\begin{equation}
    B_{t+1}=B_t\oplus\Delta B_t^{\mathrm{ev}}.
\end{equation}
Thus reliable evidence advances the local belief for every event; the later route determines whether any persistent Skill content may also change.

The two branches meet only through a typed predicate comparison,
\begin{equation}
    M_t=\mathcal{D}\left(\widehat{\Delta B}_t,
    \Delta B_t^{\mathrm{ev}},E_t\right).
\end{equation}
$M_t$ distinguishes contradiction, expected-but-unsupported, supported-but-unexpected, missing progress, termination conflict, identity conflict, temporal conflict, and uncovered predicates. In particular, unknown or insufficiently covered evidence is not negative evidence against a Skill rule.

\subsection{Hierarchical Visual Transition Credit Assignment}

The first attribution stage selects one of the three implemented routes,
{\color{red}
\begin{equation}
    z_t=A_{\mathrm{target}}(e_t,M_t)
    \in\{\mathrm{belief\mbox{-}refresh},
    \mathrm{skill\mbox{-}update},\mathrm{abstain}\}.
\end{equation}
The action interface and any environment adapter remain fixed and are not persistent update targets.}
Rule-first routing abstains when a required predicate is unknown, uncovered, or low-confidence; when a valid Skill was not followed; or when a stochastic/no-op explanation remains plausible. Identity or episode-history conflicts route to belief refresh. \red{If current goal evidence establishes completion while the pre-action belief is stale, the event likewise refreshes belief instead of modifying termination.} A constrained teacher is invoked only when these auditable rules do not resolve the event, and its answer must cite valid mismatch and evidence identifiers.

Only $z_t=\mathrm{skill\mbox{-}update}$ enters field attribution,
\begin{equation}
    h_t=A_{\mathrm{field}}(e_t,M_t)
    \in\{A,P,E,T,C\}.
\end{equation}
{\color{red}A persistent update is permitted only when the operational identifiability audit passes
\begin{equation}
    \begin{aligned}
    \mathcal{I}_{\mathrm{VTCA}}={}&
    C_{\mathrm{ev}}\wedge C_{\mathrm{prov}}\wedge C_{\mathrm{comp}}
    \wedge\neg C_{\mathrm{stoch}}\\
    &\wedge C_{\mathrm{id}}\wedge C_{\mathrm{unique}}(h_t)
    \wedge C_{\mathrm{action}}(a_t).
    \end{aligned}
\end{equation}
The seven conditions require sufficient covered evidence, complete event-local provenance, exclusion of an execution-lapse alternative, exclusion of stochasticity, resolved identity and time, exactly one implicated field, and a diagnosis bound to the current action or a compiled Skill prediction. A Skill-sourced rule may itself be refuted even when the executor followed it. Any failed condition converts the decision to abstention; this is a sufficient operational criterion under the system's observation model, not a claim of unrestricted causal identification.}

{\color{red}Skill updates carry a lifecycle label. \emph{Discovery} starts from a reliable supported-but-unexpected change: successful actions propose effect rules, whereas deterministic failed actions can expose constraint or precondition rules. Predicate arguments must generalize through the executed action arguments or the object held before the action, which prevents episode-specific identifiers and derived task-completion facts from entering persistent memory. \emph{Repair} replaces a compiled rule contradicted by independent evidence, and \emph{optimization} targets a valid but inefficient or brittle procedure.} An execution lapse remains an abstention: it adds a context-specific reminder to a decaying emphasis buffer without rewriting the canonical Skill.

\subsection{Evidence-Gated Skill Evolution}

Only events confidently attributed to the same Skill and field are eligible for Skill evolution, and repeated descriptions of one observation count as a single evidence source. \red{Events are aggregated by Skill identifier and version, attributed field, mismatch type, task pattern, and object/receptacle context; discovery events instead use an instance-free causal signature as the cluster identity. A cluster becomes ready only with unique evidence identifiers from at least two independent episodes, mean evidence confidence of at least $0.75$, and mean attribution confidence of at least $0.70$.}

\red{Each ready cluster proposes at most one bounded patch per acquisition round.} The allowed operations are append, insert after an exact statement, exact replacement, and exact deletion. A patch must cite evidence inside its cluster, touch exactly the attributed field, and match its parent version. Exact-target failure rejects the patch instead of falling back to a free-form rewrite. \red{The patch must also remain within the $512$-word active-Skill budget and contain no task-instance identifier. Discovery patches are deterministically compiled from grounded causal signatures, whereas ambiguous repair and optimization patches may use a constrained model. Textual statements, prediction rules, termination policy, and temporal rules are validated and versioned together.}

Candidate selection is staged. Static checks first validate schema, parent, field scope, provenance, and budget. Cached-transition replay then requires the candidate to repair the target event and introduce no new typed conflict. \red{The compiled candidate must also remain executable.} Surviving candidates are compared with their parent on paired proxy and disjoint finalist tasks under matched executor, seed, order, temperature, and model checkpoint. \red{For repeated rollouts, candidate--parent differences are first averaged within task and only then bootstrapped, so repeated seeds do not inflate the number of independent units.}

{\color{red}For an action-local procedure, effect, or constraint patch $p$, an outcome-independent semantic scope derived from task goals partitions paired tasks into affected set $\mathcal{D}_{A}(p)$ and protected set $\mathcal{D}_{P}(p)$ before rollout outcomes are observed. Let
\begin{equation}
    d_j(p)=\frac{1}{K}\sum_{k=1}^{K}
    \left[J_{j,k}(\omega')-J_{j,k}(\omega)\right]
\end{equation}
be the task-first paired difference across $K$ registered rollout seeds. Promotion requires
\begin{equation}
    \operatorname{LCB}_{1-\alpha}\!\left(\{d_j\}_{j\in\mathcal{D}_{A}}\right)>0,
    \qquad
    \operatorname{LCB}_{1-\alpha}\!\left(\{d_j\}_{j\in\mathcal{D}_{P}}\right)
    \geq-\epsilon.
\end{equation}
Global activation and termination patches instead require a positive global paired lower bound and no protected subgroup regression. A futility-only sequential check may reject a candidate early when an upper bound proves benefit unattainable or protected regression; it can never promote early, and every survivor must pass the full-budget bootstrap test. Pass-at-$K$ is logged only as a capability diagnostic and never enters a promotion condition.}

{\color{red}The gate records one of three dispositions. \emph{Rejected} candidates are invalid, inconsistent, incomplete, futile, or harmful. A full-budget candidate with a positive target-stream point estimate but insufficient confidence may remain \emph{shadow} only when protected and subgroup margins show no observed violation; it remains inactive and is sent to the disjoint finalist pool, where contradictory evidence rejects it. Only a candidate that passes every promotion condition becomes \emph{promoted} and replaces the active Skill. Acquisition checks ready clusters after each episode, so later evidence is collected against the newly promoted version rather than a stale batch.}

Every proposal stores its parent and candidate versions, changed field, evidence identifiers, paired metrics, and decision reason in append-only lineage. A held-out audit role evaluates proposal utility after selection without feeding results back into the gate. \red{For final evaluation, the selected Skill and executor are frozen: expected-transition construction, mismatch attribution, clustering, patching, and the evolution teacher are disabled. The agent may still update its episode-local ledger from public evidence and execute the compact textual and compiled Skill, including action-admission guards. This isolates reusable Skill improvement from test-time self-modification or extra teacher computation.}
